\section{Exact synthesis throughout the fixed-signal range}
\label{sec:entropy}
All logarithms in rates are base two. Natural logarithms are denoted
by $\ln$. Define $I$ by \eqref{eq:rate-intro}; thus $I(0)=0$, $I(1)=1$,
and $I$ is increasing on $[0,1]$.

\subsection{A uniform corner-cap lower bound}
\begin{lemma}[Entropy cap]\label{lem:entropy-cap}
For $z\in\cube^n$ and $0\le u\le1$, the fraction of sign vectors
$x$ with $x\cdot z\ge nu$ is at most $2^{-nI(u)}$. Therefore
\begin{equation}\label{eq:entropy-lower}
                         K_n(\eta)\ge2^{nI(\eta)}.
\end{equation}
\end{lemma}
\begin{proof}
For a uniform sign vector $X$ and $t\ge0$,
\[
 \E e^{tX\cdot z}=\prod_{i=1}^n\cosh(tz_i)\le\cosh(t)^n.
\]
Exponential Markov inequality gives the bound
$\exp\{-n\sup_{t\ge0}(tu-\ln\cosh t)\}$.
For $0<u<1$ the maximum occurs at $t=\operatorname{arctanh}u$.
Substitution gives $\sup(tu-\ln\cosh t)=(\ln2)I(u)$.
The two endpoints follow by limits. If a set of generators encloses
$\eta x$, then some generator $z_s$ has $x\cdot z_s\ge n\eta$.
These caps therefore cover all $2^n$ sign vectors. The fraction bound
and the union bound imply \eqref{eq:entropy-lower}.
\end{proof}
The exponent in this lemma is the classical random-access lower
exponent \cite[Theorem 2.1]{ANTVpre}. Its cap proof is useful here
because it applies directly to exact and perturbed conditional means.

\subsection{An enclosing codebook with exact, not one-sided, responses}
For $0<\eta<1$ and $0<\delta\le(1-\eta)/2$, put
\begin{align}
 k&=\left\lceil\frac{n(1+\eta+2\delta)}2\right\rceil,
 &r&=\frac{2k}n-1,\label{eq:layer-parameters}\\
 M(n,\eta,\delta)&=
 \left\lceil\frac{n+1}{\delta}
       \{n\ln(1+2/\delta)+1\}\,2^{nI(r)}\right\rceil.
                                               \label{eq:enclosure-number}
\end{align}
Here $k\le n$ and $\eta+2\delta\le r\le\eta+2\delta+2/n$.

\begin{lemma}[A Hamming-layer enclosure]\label{lem:enclosure}
There is a set of at most $M(n,\eta,\delta)$ sign vertices whose
convex hull contains $[-\eta,\eta]^n$. In particular,
\begin{equation}\label{eq:checkpoint-bound}
 K_n(\eta)\le\min\{2^n,M(n,\eta,\delta)\}.
\end{equation}
For every fixed $0<\eta<1$ there is a constant $C_\eta$ such that
\begin{equation}\label{eq:checkpoint-rate}
 nI(\eta)\le\log_2 K_n(\eta)
           \le nI(\eta)+C_\eta\log_2(n+1)\qquad(n\ge1).
\end{equation}
\end{lemma}
\begin{proof}
Fix $a\in\R^n$ with $\norm a_1=1$. Orient each coordinate by the
sign of $a_i$, choosing either orientation at zeros. Condition a uniform
sign vertex $V$ on having exactly $k$ coordinates aligned with these
orientations. Under this conditional law $Z=a\cdot V\le1$ and
$\E Z=r$. If $p=\Prb(Z\ge\eta+\delta\mid k)$, then
\[
 r\le p+(1-p)(\eta+\delta),\qquad
 p\ge\frac{r-\eta-\delta}{1-\eta-\delta}\ge\delta.
\]
It follows that, for an unconditional uniform sign vertex,
\begin{equation}\label{eq:layer-probability}
 \Prb(a\cdot V\ge\eta+\delta)
             \ge\delta\binom nk2^{-n}
             \ge\frac{\delta}{n+1}2^{-nI(r)}=:q.
\end{equation}
For the second inequality, $k$ is a mode of a binomial distribution
with success probability $k/n$, so its probability is at least
$1/(n+1)$. Rewriting that probability gives
$\binom nk\ge2^{nh_2(k/n)}/(n+1)$. The endpoints $k=0,n$ have
the same interpretation by continuity.

Choose a $\delta$-net of the $\ell^1$ unit sphere of size at most
$(1+2/\delta)^n$. Such a net follows by taking a maximal separated
set and comparing the volumes of disjoint $\ell^1$ balls of radius
$\delta/2$ with the ball of radius $1+\delta/2$.
Draw $M=M(n,\eta,\delta)$ independent uniform sign vertices. For
any net direction the probability that none has scalar product at
least $\eta+\delta$ is at most $e^{-Mq}$. The probability that
some net direction fails is at most
$(1+2/\delta)^n e^{-Mq}\le e^{-1}<1$.
Hence a fixed choice of vertices succeeds on every net direction.
Their support function is one-Lipschitz in the $\ell^1$ norm, so it
is at least $\eta$ on the entire unit sphere. Homogeneity and convex
separation imply that their hull contains the inner cube. Repeated
vertices can be discarded. Full retention gives the other upper bound.

For fixed $\eta$, take $\delta=1/n$ once $n\ge2/(1-\eta)$.
Then $r\le\eta+4/n$, and the derivative
$I'(u)=\frac12\log_2((1+u)/(1-u))$ is bounded on a fixed
neighborhood of $\eta$ lying below one. Thus
$nI(r)\le nI(\eta)+O_\eta(1)$ for sufficiently large $n$.
The remaining factor in \eqref{eq:enclosure-number} is polynomial
up to a logarithm and contributes $O(\log(n+1))$ bits.
Enlarging $C_\eta$ handles the finitely many smaller $n$.
Combine this with Lemma~\ref{lem:entropy-cap}.
\end{proof}
The hull containment gives a convex representation of \emph{each}
$\eta x$, and hence exactly the means in \eqref{eq:query-intro}.
It does not merely ensure that a coordinate is correct with probability
at least $(1+\eta)/2$. Once a successful vertex set has been chosen,
all its probabilities are fixed in the program. No random choice of a
codebook survives at runtime. For rational $\eta$, the enclosure
linear systems admit rational encoder coefficients; for algebraic
$\eta$ they admit coefficients in the corresponding ordered field.
The proof is an existence and finite-synthesis result, not a claim of
fast codebook construction.

\subsection{From a terminal codebook to a one-pass machine}
\begin{lemma}[Charged block realization]\label{lem:blocks}
Suppose the $\ell$-coordinate channel at signal $\eta$ has a
checkpoint codebook with $M$ generators. Then, for $1\le\ell\le n$,
\begin{equation}\label{eq:block-peak}
 W_n^{\rm fo}(\eta)\le
        \max\{2,\,2^\ell M^{\lfloor n/\ell\rfloor}\}.
\end{equation}
All conditional output probabilities are realized exactly.
\end{lemma}
\begin{proof}
Split the fixed input order into $b=\lfloor n/\ell\rfloor$ full
blocks and a final remainder of length less than $\ell$.
While acquiring a block, retain the previously chosen block labels and
its current raw prefix. Upon completing a block, draw one label using
that block's exact encoder, overwrite its raw buffer, and retain only
the new label with the previous labels. The encoder table is a fixed
row of the declared program. A remainder is retained raw.

At a cut the state is a tuple of at most $b$ codebook labels and at
most $\ell$ raw bits. During the overwrite the old buffer is replaced
atomically, not retained alongside a new private selector. Thus the
number of possible labels is bounded as in \eqref{eq:block-peak}.
The compression epochs are prescribed clock phases; they do not depend
on observations.

For a query in a complete block, use the corresponding label's decoder
coordinate. Its conditional mean is $\eta x_j$ independently of the
other labels. For a query in the remainder, draw a binary response of
mean $\eta x_j$ from the retained bit. At most two terminal labels
are required. The label sampling uses fresh randomness, and all its
persistent effect is in the tuple. No unread bit, chosen codebook, or
previous random seed is a free register.
\end{proof}

\begin{theorem}[Exact fixed-signal streaming rate]\label{thm:rate}
For every fixed $0<\eta<1$, there exists $C'_\eta<\infty$ such that
\begin{align}
 nI(\eta)&\le\log_2 K_n(\eta)
       \le\log_2 W_n^{\rm ad}(\eta)
       \le\log_2 W_n^{\rm fo}(\eta),\label{eq:rate-lower}\\
 \log_2 W_n^{\rm fo}(\eta)&\le
          nI(\eta)+C'_\eta\sqrt{n\log_2(n+1)}.
                                             \label{eq:rate-upper}
\end{align}
Consequently all three normalized logarithmic complexities converge
to $I(\eta)$, as in \eqref{eq:headline-rate}. The checkpoint has the
stronger error bound \eqref{eq:checkpoint-rate}.
\end{theorem}
\begin{proof}
The lower bounds have already been proved and survive adaptive
acquisition. For sufficiently large $n$, choose
$\ell=\lceil\sqrt{n\log_2(n+1)}\rceil\le n$.
Lemma~\ref{lem:enclosure} supplies a block codebook with
$\log_2 M\le\ell I(\eta)+C_\eta\log_2(\ell+1)$.
Lemma~\ref{lem:blocks} gives, after increasing an absolute constant,
\[
 \log_2 W_n^{\rm fo}
 \le\ell+\frac n\ell\{\ell I(\eta)+C_\eta\log_2(\ell+1)\}
 \le nI(\eta)+O_\eta(\sqrt{n\log(n+1)}).
\]
Again the finitely many smaller $n$ are absorbed into $C'_\eta$.
Division by $n$ proves the limits.
\end{proof}
This theorem covers every fixed interior signal, rather than just two
endpoint windows. It does not determine every finite integer value of
$K_n$ or $W_n$, their lower-order difference, or a uniform sharp law
when $\eta$ varies with $n$. Equations~\eqref{eq:enclosure-number} and
\eqref{eq:block-peak} remain explicit nonasymptotic bounds in those
varying-signal problems. In the strong finite regime
$\eta>1-1/n$, the disjoint-cap argument gives $K_n=W_n^{\rm ad}
=W_n^{\rm fo}=2^n$: the encoder average of
$\frac12\sum_i(1-x_i z_i)$ is below $1/2$, and its strict caps for
distinct input corners are disjoint. Full retention matches this bound.
That sharper endpoint statement is consistent with, but not needed for,
Theorem~\ref{thm:rate}.
